\section{Failure analysis}
The two lowest-scoring photos in Table~\ref{tab:cap} are the two failure cases (Figures~\ref{fig:ov1} and~\ref{fig:ov2}).
\begin{itemize}
\item \textbf{Clutter and many similar objects} (\texttt{cycle\_inside}, $Q=0.333$, BF1 $0.078$). The ground truth marks only the large bicycle in the right foreground, but the shed contains a whole row of bicycles. The network instead fires on the dense group of bicycles in the background and largely misses the foreground one. The task has no notion of ``the object the user means''; DUTS teaches saliency as a high-contrast, centred object, and with many equal candidates a small network with little global context picks the textured cluster. The same pattern lowers \texttt{study\_table} ($Q=0.369$) and \texttt{chair\_study-table}.
\item \textbf{Low-contrast, irregular object} (\texttt{outside\_plant}, $Q=0.365$, BF1 $0.108$). The ground truth is the row of shrubs, which resemble the soil and weeds in colour and texture. The prediction misses the shrubs and instead selects the bright concrete footpath. Foliage boundaries are thin and ragged, so even where overlap is acceptable the contour score is very low.
\end{itemize}
A milder issue is bleeding onto the supporting surface or neighbouring structure (\texttt{shoes}, \texttt{dustbin}: a small spur near the wall), probably because object and support share edges and colour statistics. The best case is the isolated, high-contrast \texttt{dustbin} ($Q=0.854$). Because the ground-truth masks are AI-generated references, part of the remaining error in any photo may come from the reference rather than the model.
